\chapter{आकार और स्थान}\label{ch:g1:shapes}

खिड़कियाँ आयत हैं, पहिए वृत्त हैं, छतों में त्रिभुज छिपे हैं:
आकार हर जगह हैं। इस अध्याय में हम भुजाएँ और कोने गिनकर उन्हें
पहचानना सीखेंगे, और ठीक-ठीक कहना सीखेंगे कि चीज़ें \emph{कहाँ}
हैं: बाएँ, दाएँ, ऊपर, नीचे, बीच में।

\section{चार बुनियादी आकार}

\begin{definition}[वर्ग, आयत, त्रिभुज, वृत्त]\label{def:g1:shapes:def}
\begin{itemize}
  \item \emph{त्रिभुज}\index{त्रिभुज} की $3$ भुजाएँ और $3$ कोने होते हैं;
  \item \emph{वर्ग}\index{वर्ग} की $4$ बराबर भुजाएँ और $4$ कोने होते हैं;
  \item \emph{आयत}\index{आयत} की $4$ भुजाएँ और $4$ कोने होते हैं, और
        आमने-सामने की भुजाएँ बराबर होती हैं (खिंचा हुआ वर्ग);
  \item \emph{वृत्त}\index{वृत्त} पूरी तरह गोल है: न भुजा, न कोना।
\end{itemize}
\end{definition}

\begin{omfigure}
\begin{tikzpicture}[scale=0.7]
  \draw[thick, fill=omDef!12] (0,0) -- (2,0) -- (1,1.7) -- cycle;
  \node[below, align=center, font=\small] at (1,-0.25)
    {त्रिभुज:\\$3$ भुजाएँ};
  \begin{scope}[xshift=3.6cm]
  \draw[thick, fill=omThm!12] (0,0) rectangle (1.7,1.7);
  \node[below, align=center, font=\small] at (0.85,-0.25)
    {वर्ग:\\$4$ बराबर भुजाएँ};
  \end{scope}
  \begin{scope}[xshift=6.9cm]
  \draw[thick, fill=omMeth!12] (0,0.05) rectangle (2.8,1.65);
  \node[below, align=center, font=\small] at (1.4,-0.25)
    {आयत:\\$4$ भुजाएँ};
  \end{scope}
  \begin{scope}[xshift=11.3cm]
  \draw[thick, fill=omProp!12] (0.9,0.85) circle (0.9);
  \node[below, align=center, font=\small] at (0.9,-0.25)
    {वृत्त:\\कोई कोना नहीं};
  \end{scope}
\end{tikzpicture}

{\small चार बुनियादी आकार। आकार का नाम बताने के लिए पहली नज़र पर
भरोसा मत करो: उसकी भुजाएँ और कोने \emph{गिनो}।}
\end{omfigure}

\begin{example}[तिरछा आकार भी वही आकार है]\label{ex:g1:shapes:tilted}
कोने पर खड़ा वर्ग भी वर्ग ही है --- चार बराबर भुजाएँ, चार
कोने --- भले ही वह पतंग जैसा दिखे। आकार को घुमाने से वह बदल
नहीं जाता।
\end{example}

\begin{example}[आकारों की खोज]\label{ex:g1:shapes:hunt}
कक्षा में: बोर्ड आयत है, घड़ी वृत्त, मुड़ा हुआ रुमाल त्रिभुज बन
सकता है, कुछ खिड़कियाँ वर्ग हैं। असली चीज़ें पूरी तरह सटीक आकार
नहीं होतीं, पर उनके फलक इनके क़रीब होते हैं (ठोस और उनके फलक
\cref{ch:g2:shapes} में लौटेंगे)।
\end{example}

\section{यह कहाँ है?}

\begin{definition}[स्थान बताने वाले शब्द]\label{def:g1:shapes:position}
कोई चीज़ कहाँ है, यह बताने के लिए: \emph{बाएँ} / \emph{दाएँ},
\emph{ऊपर} / \emph{नीचे}, \emph{पर} / \emph{तले},
\emph{सामने} / \emph{पीछे}, \emph{बीच में},
\emph{अंदर} / \emph{बाहर}।
\end{definition}

\begin{omfigure}
\begin{tikzpicture}[scale=0.75]
  \draw[thick, fill=omDef!12] (0,0) circle (0.55);
  \draw[thick, fill=omThm!12] (1.9,-0.55) rectangle (3.0,0.55);
  \draw[thick, fill=omMeth!12] (4.4,-0.5) -- (5.5,-0.5) -- (4.95,0.45)
    -- cycle;
  \node[below, font=\small] at (0,-0.75) {वृत्त};
  \node[below, font=\small] at (2.45,-0.75) {वर्ग};
  \node[below, font=\small] at (4.95,-0.75) {त्रिभुज};
  \node[right, font=\small, align=left] at (6.3,0)
    {वर्ग, वृत्त और त्रिभुज के\\बीच में है;\\
    वृत्त वर्ग के बाईं ओर है};
\end{tikzpicture}

{\small स्थान बताने वाले शब्द काम पर। ध्यान रहे: \emph{तुम्हारा}
बायाँ और तुम्हारे सामने खड़े किसी का बायाँ उलटे होते हैं!}
\end{omfigure}

\section{जाल पर रास्ते}

\begin{method}[रास्ता बताना]\label{met:g1:shapes:path}
जाल वाले काग़ज़ पर रास्ता कदमों की एक सूची है:
$\to$ (एक खाना दाईं ओर), $\leftarrow$ (बाईं ओर), $\uparrow$ (ऊपर),
$\downarrow$ (नीचे)।
\begin{enumerate}
  \item रास्ता चलने के लिए: निशान लगे खाने से शुरू करो और तीरों
        का एक-एक करके पालन करो;
  \item रास्ता बताने के लिए: मन-ही-मन चलो और तीरों को क्रम से
        लिखो;
  \item दो अलग-अलग रास्ते एक ही दो खानों को जोड़ सकते हैं ---
        कोई छोटा, कोई लंबा।
\end{enumerate}
\end{method}

\begin{omfigure}
\begin{tikzpicture}[scale=0.6]
  \draw[gray!40] (0,0) grid (7,4);
  \fill[omDef!40] (0.5,0.5) +(-0.42,-0.42) rectangle +(0.42,0.42);
  \node[font=\footnotesize] at (0.5,0.5) {शुरू};
  \fill[omThm!40] (5.5,2.5) +(-0.42,-0.42) rectangle +(0.42,0.42);
  \node[font=\footnotesize] at (5.5,2.5) {लक्ष्य};
  % सात इकाई तीर, हर कदम के लिए एक, बीच में छोटी ख़ाली जगह
  \foreach \a/\b in {0.5/1.5, 1.5/2.5}
    \draw[-Stealth, omDef, very thick] (\a+0.1,0.5) -- (\b-0.1,0.5);
  \foreach \a/\b in {0.5/1.5, 1.5/2.5}
    \draw[-Stealth, omDef, very thick] (2.5,\a+0.1) -- (2.5,\b-0.1);
  \foreach \a/\b in {2.5/3.5, 3.5/4.5}
    \draw[-Stealth, omDef, very thick] (\a+0.1,2.5) -- (\b-0.1,2.5);
  \draw[-Stealth, omDef, very thick] (4.6,2.5) -- (5.05,2.5);
\end{tikzpicture}

{\small शुरू से लक्ष्य तक एक रास्ता, हर कदम पर एक तीर:
$\to\ \to\ \uparrow\ \uparrow\ \to\ \to\ \to$ --- सात कदम। क्या
तुम इतने ही कदमों वाला दूसरा रास्ता खोज सकते हो?}
\end{omfigure}

\section{अभ्यास}

\begin{exercise}[$\star$]\label{exo:g1:shapes:1}
हर आकार की भुजाएँ और कोने गिनकर उसका नाम बताओ: $3$ कोनों वाला
आकार; बिना कोने का गोल आकार; $4$ बराबर भुजाओं वाला
आकार।
\end{exercise}

\begin{exercise}[$\star$]\label{exo:g1:shapes:2}
घर में खोजो: दो आयत, एक वृत्त, एक त्रिभुज, एक वर्ग।
(दरवाज़े? थालियाँ? सड़क के चिह्न?)
\end{exercise}

\begin{exercise}[$\star$]\label{exo:g1:shapes:3}
सही या ग़लत? ``कोने पर घुमाया गया वर्ग, वर्ग नहीं रह
जाता।'' \cref{ex:g1:shapes:tilted} की मदद से समझाओ।
\end{exercise}

\begin{exercise}[$\star$]\label{exo:g1:shapes:4}
जाल वाले काग़ज़ पर बनाओ: $3$ गुणा $3$ खानों का वर्ग; $5$ गुणा
$2$ का आयत; जाल की एक तिरछी रेखा को भुजा बनाकर एक त्रिभुज।
\end{exercise}

\begin{exercise}[$\star$]\label{exo:g1:shapes:5}
तीन खिलौने एक पंक्ति में रखो। स्थान बताने वाले शब्दों से पंक्ति
बताओ: बीच में कौन है? बाईं ओर कौन है?
\end{exercise}

\begin{exercise}[$\star$]\label{exo:g1:shapes:6}
एक वृत्त बनाओ। उसके \emph{अंदर} एक क्रॉस, \emph{बाहर} एक तारा,
और ठीक वृत्त \emph{पर} एक बिंदु बनाओ।
\end{exercise}

\begin{exercise}[$\star$]\label{exo:g1:shapes:7}
जाल वाले काग़ज़ पर शुरू का एक खाना चुनो। चलो:
$\to\ \uparrow\ \to\ \uparrow\ \to$। पहुँचने की जगह पर निशान लगाओ। फिर
वह रास्ता लिखो जो सीधे शुरू पर लौटा दे।
\end{exercise}

\begin{exercise}[$\star$]\label{exo:g1:shapes:8}
कुल कितने कोने: दो त्रिभुज और एक वर्ग? एक वृत्त और दो
आयत?
\end{exercise}

\begin{exercise}[$\star\star$]\label{exo:g1:shapes:9}
$12$ तीलियाँ लो। क्या \emph{सारी} तीलियों से, पूरी तीलियाँ रखते
हुए और हर भुजा पर बराबर संख्या रखते हुए, वर्ग बना सकते हो? हर
भुजा पर कितनी तीलियाँ? क्या $10$ तीलियों से यह हो सकेगा?
\end{exercise}
